\chapter{Analysen}
\label{kap:format-analysen}

\section{Übersicht}

Die Liste \code{analyses} nennt die Analysen; sie werden in dieser
Reihenfolge gerechnet.

\begin{center}
\begin{tabularx}{\textwidth}{@{}lP{5.6cm}L@{}}
  \toprule
  \code{type} & Parameter & Ergebnis \\
  \midrule
  \code{op} & – & Arbeitspunkt (Konsole) \\
  \code{dc} & \code{source}, \code{start}, \code{stop}, \code{step} & Sweep einer V- oder I-Quelle \\
  \code{ac} & \code{sweep}, \code{points}, \code{fstart}, \code{fstop} & Betrag und Phase je Ausgang \\
  \code{tran} & \code{tstop}, optional \code{tstep}, \code{method}, \code{uic}, \code{output_step}, \dots & Zeitverlauf \\
  \code{noise} & \code{output}, \code{input}, Frequenzen wie \code{ac} & Rauschdichten \\
  \code{pss} & \code{period}, optional \code{periods}, \code{tstab}, \code{tol} & eingeschwungene Periode \\
  \code{sens} & \code{analysis}, optional \code{parameters}, \code{delta} & Empfindlichkeiten \\
  \code{fit} & \code{analysis}, \code{output}, \code{data}, \code{parameters} & angepasste Parameter \\
  \bottomrule
\end{tabularx}
\end{center}

Für jede Analyse gilt außerdem:

\begin{itemize}
  \item \code{"name"} gibt ihr einen eigenen Namen (für die Ergebnisdatei
    und für Python); ohne Namen heißt sie wie ihr Typ.
  \item \code{"outputs"} ersetzt für diese Analyse die globale Liste der
    Ausgänge.
  \item \code{"measure"} berechnet Kennwerte aus den Kurven
    (Abschnitt~\ref{sec:messgroessen}).
  \item \code{"param_sweep"} und \code{"monte_carlo"} wiederholen sie
    (Kapitel~\ref{kap:format-studien}).
  \item \code{"save": false} unterdrückt die Datei mit den Kurven, wenn nur
    Messgrößen interessieren.
\end{itemize}

\section{Arbeitspunkt (op)}

\begin{json}
{ "type": "op" }
\end{json}

Der Arbeitspunkt ist die stationäre Lösung: Kondensatoren sind offen,
Spulen kurzgeschlossen, Quellen stehen auf ihrem \code{dc}-Wert,
Zeitableitungen eigener Gleichungen sind null. Die Werte der Ausgänge
erscheinen auf der Konsole. Konvergiert das Newton-Verfahren nicht direkt,
versucht esim nacheinander eine Pseudo-Transiente, gmin-Stepping und
Quellen-Stepping (Kapitel~\ref{kap:numerik}); die Konsole nennt, welches
Verfahren zum Ziel geführt hat. Bei bistabilen Schaltungen wählt
\code{"nodeset"} den gewünschten Zustand (Abschnitt~\ref{sec:optionen}).

\section{DC-Sweep (dc)}

\begin{json}
{ "type": "dc", "source": "Vin", "start": 0, "stop": 5, "step": 0.01 }
\end{json}

Der Wert der Quelle \code{source} (Spannungs- oder Stromquelle, auch in
einer Teilschaltung: \code{"X1.V1"}) läuft von \code{start} bis \code{stop}
in Schritten von \code{step}; für jeden Wert wird der Arbeitspunkt
bestimmt, ausgehend vom vorigen. Die erste Spalte der Ergebnisdatei ist der
Wert der Quelle. Übertragungskennlinien, Ausgangskennlinienfelder (mit einem
Parameter-Sweep außen herum) und Kennlinien von Bauteilen entstehen so.

\section{Frequenzgang (ac)}

\begin{json}
{ "type": "ac", "sweep": "dec", "points": 20, "fstart": 10, "fstop": "10meg" }
\end{json}

Die Kleinsignalanalyse linearisiert die Schaltung im Arbeitspunkt und
rechnet für jede Frequenz das komplexe lineare Gleichungssystem. Angeregt
wird sie von den Quellen mit \code{ac} (und \code{acphase}).

\begin{itemize}
  \item \code{sweep}: \code{"dec"} (Standard; \code{points} je Dekade),
    \code{"oct"} (je Oktave) oder \code{"lin"} (\code{points} insgesamt).
    \code{points} ist ohne Angabe 10.
  \item Die Ergebnisdatei enthält je Ausgang \code{mag(...)} (Betrag) und
    \code{phase(...)} (Grad). Der Editor zeigt den Betrag in dB.
  \item Nur einfache Ausgänge sind erlaubt: \code{v(knoten)},
    \code{v(a,b)}, \code{i(bauteil)}, Variablen.
  \item Eigene Gleichungen und Verhaltensquellen sind in der
    Linearisierung enthalten. Die Laufzeit von Leitungen wirkt exakt als
    Phasendrehung; \code{delay(x, td)} in Ausdrücken ist in der
    AC-Analyse dagegen gleich \code{x}.
\end{itemize}

\section{Zeitverlauf (tran)}
\label{sec:tran}

\begin{json}
{ "type": "tran", "tstop": "5m", "tstep": "10u" }
\end{json}

\begin{longtable}{@{}lP{10.6cm}@{}}
  \toprule
  Schlüssel & Bedeutung \\
  \midrule
  \endhead
  \bottomrule
  \endfoot
  \code{tstop} & Endzeit. \\
  \code{tstep} & Gewünschte Auflösung: Kein Schritt wird länger (Standard \code{tstop/50}). \\
  \code{hmax} & Größte Schrittweite (wie \code{tstep}, nur ohne Bezug zur Ausgabe). \\
  \code{hinit}, \code{hmin} & Erste und kleinste Schrittweite. \\
  \code{method} & Integrationsverfahren: \code{"bdf"} (Standard), \code{"gear2"}, \code{"trap"}, \code{"be"}. \\
  \code{max_order} & Höchste Ordnung von \code{"bdf"} (1 bis 5, Standard 5). \\
  \code{uic} & \code{true}: ohne Arbeitspunkt von den Anfangsbedingungen starten. \\
  \code{output_step} & Ausgabe auf dem festen Raster 0, $\Delta t$, $2\Delta t$, \dots \\
  \code{max_events} & Höchstzahl von Ereignissen (Standard 100\,000). \\
  \code{fourier} & Fourier-Analyse am Ende (Abschnitt~\ref{sec:fourier}). \\
\end{longtable}

\paragraph{Start.}
Die Schaltung beginnt in ihrem Arbeitspunkt (mit den Werten der Quellen bei
$t=0$), Zustände eigener Gleichungen bei ihren Anfangswerten
(Abschnitt~\ref{sec:anfangswerte}). Mit \code{"uic": true} wird die
Anfangslösung übersprungen: Es gelten \code{ic} an Kondensatoren und
Spulen, \code{"ic": {"knoten": spannung}} auf oberster Ebene und die
Startwerte der Variablen. Knoten, die Spannungsquellen von Masse aus
festlegen (Versorgungen, auch über Strommesser hinweg), starten auf deren
Spannung; alles andere bei 0.

\paragraph{Schrittweite.}
Die Schrittweite wird automatisch aus dem geschätzten Fehler gewählt
(Abschnitt~\ref{sec:genauigkeit}). Ecken der Quellen (Pulsflanken,
PWL-Stützstellen), Ereignisse und Laufzeiten von Leitungen und Totzeiten
werden exakt getroffen. Sinusquellen begrenzen den Schritt auf ein
Zwanzigstel ihrer Periode.

\paragraph{Verfahren.}
\begin{description}[style=nextline,leftmargin=2em]
  \item[\code{bdf} (Standard)]
    BDF mit variabler Ordnung 1 bis 5. Nach jedem Schritt wird die
    Ordnung gewählt, die den größten nächsten Schritt erlaubt. Bei glatten
    Lösungen und engen Toleranzen braucht es ein Vielfaches weniger
    Schritte als ein Verfahren fester Ordnung.
  \item[{\code{gear2}}]
    BDF fester Ordnung 2. Robust, klingelt nicht, dämpft Schwingungen
    leicht.
  \item[{\code{trap}}]
    Trapezregel. Keine numerische Dämpfung – die beste Wahl für schwach
    gedämpfte Schwingkreise über viele Perioden. Dafür können Ströme von
    Spannungsquellen von Schritt zu Schritt leicht alternieren.
  \item[{\code{be}}]
    Rückwärts-Euler: erste Ordnung, sehr robust, stark dämpfend.
\end{description}

\paragraph{Ausgabe.}
Ohne \code{output_step} wird jeder Zeitpunkt des Lösers ausgegeben; die
Abstände sind dann ungleich. Mit \code{output_step} wird auf ein festes
Raster interpoliert, was für die Weiterverarbeitung (FFT, Vergleich mit
Messungen) praktisch ist; die Schrittweite des Lösers bleibt dabei frei.
Am Zeitpunkt eines Ereignisses steht der Zustand \emph{vor} der Zuweisung,
im Punkt unmittelbar danach der neue.

\section{Fourier-Analyse}
\label{sec:fourier}

\begin{json}
{ "type": "tran", "tstop": "5m", "tstep": "10u",
  "fourier": { "freq": "1k", "harmonics": 9, "outputs": ["v(out)"] } }
\end{json}

Der Block \code{fourier} einer Transientenanalyse berechnet das Spektrum
der letzten ganzen Periode(n) zur Grundfrequenz \code{freq}. Das Ergebnis
steht in \datei{<Name>.tran.fourier.csv} (Betrag und Phase je Harmonischer,
Zeile 0 ist der Mittelwert); Grundschwingung und Klirrfaktor erscheinen auf
der Konsole. \code{periods} (Standard 1) wählt die Zahl der ausgewerteten
Perioden, \code{from} den frühesten Zeitpunkt. Die Genauigkeit hängt von der
Auflösung des Zeitverlaufs ab: Für kleine Klirrfaktoren \code{tstep} und
\code{reltol} klein wählen.

\section{Rauschen (noise)}

\begin{json}
{ "type": "noise", "output": "v(out)", "input": "Vin",
  "sweep": "dec", "points": 20, "fstart": 10, "fstop": "100k" }
\end{json}

Berechnet im Arbeitspunkt die spektrale Rauschdichte am Ausgang
(\code{onoise}, bei einer Spannung in V/$\sqrt{\text{Hz}}$) und, wenn eine
Quelle als \code{input} genannt ist, die auf sie bezogene Dichte
(\code{inoise}) samt Verstärkung. Auf der Konsole stehen das über das Band
integrierte Rauschen und die größten Beiträge. \code{output} ist
\code{v(knoten)}, \code{v(a,b)} oder \code{i(bauteil)}.

Rauschquellen sind das thermische Rauschen aller Widerstände (auch der
Bahnwiderstände von Diode und Transistor), das Schrotrauschen von
Dioden-, Kollektor- und Basisstrom, das thermische Kanalrauschen des MOSFET
und 1/f-Rauschen über die Modellparameter \code{kf} und \code{af}.
Schalter, ideale Dioden und Verhaltensquellen rauschen nicht; Widerstände
mit \code{"noise": false} auch nicht.

\section{Periodischer stationärer Zustand (pss)}

Schaltnetzteile, Gleichrichter oder Verstärker mit großen
Koppelkondensatoren brauchen oft Tausende Perioden, bis sie eingeschwungen
sind. Die Analyse \code{pss} sucht den eingeschwungenen Zustand direkt:

\begin{json}
{ "type": "pss", "period": "10u",
  "measure": { "vout": { "avg": "v(out)" }, "ripple": { "pp": "v(out)" } } }
\end{json}

\begin{itemize}
  \item \code{period} ist die Periode der Anregung. Alle Quellen müssen sich
    mit ihr wiederholen; das wird geprüft.
  \item Gesucht wird ein Zustand, der nach einer Periode wieder erreicht
    ist. Jede Iteration rechnet eine Periode und schätzt aus den letzten
    Ergebnissen, wohin die langsamen Vorgänge laufen
    (Anderson-Beschleunigung). Der Tiefsetzsteller aus
    \datei{examples/buck_converter.json} braucht so 5 Perioden statt etwa
    2000.
  \item Ergebnis ist eine Periode im eingeschwungenen Zustand
    (\code{"periods": n} für mehrere); Messgrößen, \code{fourier} und
    \code{output_step} wirken wie bei \code{tran}.
  \item \code{tol} (Standard $10^{-4}$) ist die relative Abweichung, mit
    der sich der Zustand nach einer Periode wiederholen muss;
    \code{max_iterations} begrenzt die Zahl der Iterationen. \code{tstab}
    rechnet vor der Suche eine Zeit lang normal, was bei stark
    nichtlinearem Anlauf hilft.
\end{itemize}

Die Periode muss bekannt sein. Für selbstschwingende Schaltungen
(Oszillatoren) ist die Analyse nicht gedacht, ebenso wenig für Zustände, die
erst nach einem Schaltvorgang mit Gedächtnis erreicht werden (Hysterese,
diskrete Variablen). Totzeiten und Leitungen unterstützt sie nicht.

\section{Empfindlichkeiten (sens)}

Die Analyse \code{sens} beantwortet die Frage, welche Parameter ein
Ergebnis bestimmen:

\begin{json}
{ "type": "sens", "parameters": ["R1", "R2", "C"],
  "analysis": { "type": "ac", "sweep": "dec", "points": 100, "fstart": 1, "fstop": "100k",
                "measure": { "fc": { "f3db": "v(out)" } } } }
\end{json}

\begin{itemize}
  \item Jeder Parameter wird um $\pm$\code{delta} (relativ, Standard
    $10^{-3}$) geändert, die eingebettete Analyse wiederholt und der
    Differenzenquotient gebildet – zwei Rechnungen je Parameter.
  \item Betrachtet werden die Messgrößen der Analyse, bei \code{op} ihre
    Ausgänge.
  \item Ohne \code{parameters} werden alle Einträge von \code{parameters}
    variiert.
  \item Ausgegeben wird die normierte Empfindlichkeit: um wie viel Prozent
    sich das Ergebnis ändert, wenn der Parameter um ein Prozent wächst. Die
    Tabelle steht in \datei{<Name>.sens.csv}; die Zusammenfassung
    (\code{--summary}) enthält zusätzlich die absoluten Ableitungen.
  \item Messgrößen aus abgetasteten Kurven (Grenzfrequenzen,
    Durchgangszeiten) sind nur so genau wie das Raster; dann ein größeres
    \code{delta} (z.\,B. 0,02) oder ein feineres Raster wählen.
\end{itemize}

\section{Anpassung an Messdaten (fit)}

Die Analyse \code{fit} verändert Parameter so, dass eine simulierte Kurve
möglichst gut auf Datenpunkten liegt (kleinste Quadrate nach
Levenberg-Marquardt):

\begin{json}
{ "type": "fit",
  "analysis": { "type": "dc", "source": "V1", "start": 0.25, "stop": 0.75, "step": 0.005 },
  "output": "-i(V1)",
  "data": { "file": "data/diode_iv.csv", "xcol": 0, "ycol": 1 },
  "parameters": { "is0": { "min": 1e-14, "max": 1e-6, "log": true },
                  "n0":  { "min": 1, "max": 3 } },
  "residual": "log" }
\end{json}

\begin{itemize}
  \item \code{analysis} ist eine \code{dc}-, \code{ac}- oder
    \code{tran}-Analyse; \code{output} die Größe, die mit den Daten
    verglichen wird (bei \code{ac} ihr Betrag). Die Daten müssen innerhalb
    des simulierten Bereichs liegen; dazwischen wird linear interpoliert.
  \item \code{data}: \code{x} und \code{y}, \code{points} oder eine Datei
    wie bei den Tabellen (Abschnitt~\ref{sec:tabellen}).
  \item \code{parameters}: je Parameter optional \code{start} (sonst der
    Wert aus \code{parameters}), \code{min}, \code{max} und
    \code{"log": true}, wenn sich der Wert über Größenordnungen bewegen
    kann.
  \item \code{residual}: \code{"lin"} (Differenz, Standard), \code{"rel"}
    (relativ zu den Daten) oder \code{"log"} (Logarithmus des Verhältnisses,
    für Daten über mehrere Dekaden).
  \item Ergebnis: die angepassten Werte und der mittlere Restfehler auf der
    Konsole, Daten und Modellkurve in \datei{<Name>.fit.csv}.
\end{itemize}

Das Verfahren findet ein lokales Minimum; die Startwerte sollten in der
richtigen Größenordnung liegen. Parameter, die die Kurve im gemessenen
Bereich kaum beeinflussen, bleiben unbestimmt.
\datei{examples/fit_diode.json} passt so das Modell einer Diode an eine
gemessene Kennlinie an.

\section{Messgrößen}
\label{sec:messgroessen}

Mit \code{measure} berechnet eine Analyse Kennwerte aus ihren Kurven:

\begin{json}
{ "type": "tran", "tstop": "200m",
  "measure": {
    "vdc":    { "avg": "v(out)", "from": "160m" },
    "ripple": { "pp":  "v(out)", "from": "160m" },
    "t90":    { "cross": "v(out)", "level": 13.5 }
  } }
\end{json}

Jede Messung besteht aus einer Funktion mit dem auszuwertenden Ausdruck und
Optionen. $x$ ist dabei die Zeit, die Frequenz oder der Wert der gesweepten
Quelle; in der AC-Analyse wird der Betrag ausgewertet.

\begin{longtable}{@{}P{3.2cm}P{6.4cm}P{4.2cm}@{}}
  \toprule
  Funktion & Ergebnis & Optionen \\
  \midrule
  \endhead
  \bottomrule
  \endfoot
  \code{min}, \code{max}, \code{pp} & Extremwerte, Spitze-Spitze & \\
  \code{argmin}, \code{argmax} & $x$, bei dem der Extremwert liegt & \\
  \code{avg}, \code{rms}, \code{integ} & Mittelwert, Effektivwert, Integral über $x$ & \\
  \code{final} & letzter Wert & \\
  \code{at} & Wert an der Stelle \code{x} & \code{x} \\
  \code{cross} & $x$ des \code{n}-ten Durchgangs durch \code{level} & \code{level}, \code{n} (Standard 1), \code{edge}: \code{rise}, \code{fall}, \code{both} \\
  \code{period}, \code{freq}, \code{duty} & Periode, Frequenz, Tastgrad aus den steigenden Durchgängen & \code{level} (Standard: Mitte zwischen Minimum und Maximum) \\
  \code{rise}, \code{fall} & Anstiegs- bzw. Abfallzeit zwischen zwei Pegeln & \code{low}, \code{high} (Standard 10\,\% und 90\,\% des Hubs) \\
  \code{f3db}, \code{f3db_low} & obere bzw. untere Frequenz, bei der die Kurve auf Maximum$/\sqrt2$ gefallen ist & \\
  \code{thd}, \code{harmonic} & Klirrfaktor (bis zur 15. Harmonischen) bzw. Amplitude der \code{n}-ten Harmonischen & \code{fundamental}, \code{n} \\
\end{longtable}

\code{from} und \code{to} beschränken jede Messung auf einen Ausschnitt von
$x$. Eine Messung ohne Ergebnis (z.\,B. kein Durchgang gefunden) liefert
»not defined« bzw. \code{nan}. Ausgewertet werden die ausgegebenen Punkte,
mit \code{output_step} also das Raster.

Ohne Studie erscheinen die Messgrößen auf der Konsole und im Editor in der
Wertetabelle. In Parameter-Sweeps und Monte-Carlo-Läufen entsteht daraus
eine Tabelle mit einer Zeile je Lauf (Kapitel~\ref{kap:format-studien}).
